\documentclass{article}

\usepackage[margin=1in]{geometry}
\usepackage{xcolor}
\usepackage{stepdiff}

\title{stepdiff Relation Alignment Examples}
\author{}
\date{}

\stepdiffsetup{display=notes, theme=soft, reason-style=muted}

\begin{document}

\maketitle

\section*{Inequalities}

\[
\begin{stepdiff}[frame=true]
\step{a_n \le b_n}
\step[reason={take limits}]{\lim a_n \le \lim b_n}
\step[reason={conclude}, tag=final]{L \le M}
\end{stepdiff}
\]

\[
\begin{stepdiff}[display=focus, theme=focus]
\step{x > 0}
\step[reason={square both sides}]{x^2 > 0}
\step[reason={weaken}, tag=final]{x^2 \ge 0}
\end{stepdiff}
\]

\section*{Approximations and equivalences}

\[
\begin{stepdiff}[highlight=color]
\step{\sin x \sim x}
\step[reason={small angle}]{\sin x \approx x}
\step[reason={at zero}, tag=final]{\sin 0 = 0}
\end{stepdiff}
\]

\[
\begin{stepdiff}[theme=minimal, highlight=underline]
\step{n \equiv 1 \pmod 2}
\step[reason={write as odd}, tag=final]{n = 2k+1}
\end{stepdiff}
\]

\section*{Implications}

\[
\begin{stepdiff}[display=slide, theme=focus, reason-style=badge]
\step{P \Rightarrow Q}
\step[reason={add hypothesis}]{P \wedge R \Rightarrow Q}
\step[reason={chain implication}, tag=final]{P \wedge R \Longrightarrow S}
\end{stepdiff}
\]

\end{document}
